\begin{annexitem}{annex:p19.rat132.order}{thm:p19-rat132-order}{Order consequence without causal import}
\begin{proof}
Transitivity of the strict order gives \(c<g\) from \(c<m<g\).
For the fixed northward coordinate interpretation, if
\(y_A>y_B\) and \(y_B>y_C\), transitivity gives \(y_A>y_C\).
This uses the same coordinate frame in both premises.

To test the causal conclusion, expand one ordered-event structure with an
additional causal-influence relation. In one expansion put \((c,g)\)
in it; in another let that relation be empty. Both expansions retain the
identical strict event order and satisfy both premises, but the proposed
causal assertion is false in the second. Hence the order premises do not
entail it. No claim that the actual historical events were causally
independent is needed for this nonentailment argument.
\end{proof}
\end{annexitem}

\begin{annexitem}{annex:p19.rat132.endpoints}{thm:p19-rat132-endpoints}{Thirteen-way partition and a complete finite search}
\begin{proof}
Compare \(i^-\) with \(j^-\). If \(i^-<j^-\), compare
\(i^+\) with \(j^-\): less gives before, equality meets, and greater
leaves a comparison with \(j^+\). In the latter case less gives overlaps,
equality finished-by, and greater contains. If \(i^-=j^-\), comparison
of the ends gives starts, equals or started-by. If \(j^-<i^-\),
apply the first case with exchanged intervals, giving after, met-by,
overlapped-by, finishes, or during. The branches are disjoint and cover
all comparisons; positive lengths exclude other cases. These are six
relations, their six inverses and equality.

For \(n\) intervals, there are \(2n\) named endpoints. Enumerate all
functions from these endpoint names to \(\{1,\ldots,2n\}\). Discard
functions violating any interval's strict start-before-end condition;
interpret comparisons of integer values as endpoint comparisons and test
every relation-set constraint by the table. This finite enumeration may
represent the same weak order more than once; duplication does not affect
termination or correctness. Every retained satisfying assignment gives
proper real intervals and so is a sound witness.

Conversely, any real solution has finitely many distinct endpoint values.
Sort them increasingly and assign ranks \(1,\ldots,k\), where
\(k\leq2n\). Equal real endpoints receive the same rank. This preserves
every strict comparison and equality, hence every base relation and each
disjunction of them. The rank assignment occurs in the enumeration, so
no solution is missed. For zero intervals the empty assignment supplies
the same conclusion directly.

Consistency holds exactly when the surviving assignments are nonempty.
Possibility of \(q\) holds exactly when some surviving assignment satisfies
\(q\). A necessary query on a separately verified satisfiable premise set
holds exactly when all surviving assignments satisfy \(q\); equivalently,
there is no survivor satisfying its negation. For a base-relation query,
negation is the complementary set among the thirteen relations. The
finite language's endpoint semantics justifies these tests; it does not
justify ignoring endpoint equality or generalizing to an unprovided RCC
language.
\end{proof}
\end{annexitem}

\begin{annexitem}{annex:p19.rat132.propagation}{thm:p19-rat132-propagation}{Soundness and termination of local interval narrowing}
\begin{proof}
Fix a satisfying assignment before an update. Its actual relation between
\(i,j\) lies in \(N_{ij}\), and that between \(j,k\) lies in
\(N_{jk}\). The three assigned intervals themselves witness that their
actual \(i,k\) relation belongs to the composition
\(N_{ij}\diamond N_{jk}\). It also lies in the current \(N_{ik}\)
by satisfaction. Thus it lies in their intersection after the update.
The converse update merely exchanges the two assigned intervals and is
also sound. Intersecting with a new asserted premise preserves precisely
the assignments satisfying that added premise; it is not asserted to
preserve assignments violating new information.

Every strict change deletes at least one of the thirteen base relations
from one of finitely many arcs. At most \(13n(n-1)\) off-diagonal
deletions are available, and diagonal equality is fixed. A finite worklist
enqueued only on strict change therefore terminates; no-op comparisons do
not re-enqueue themselves. If any arc is empty no actual interval relation
can satisfy it, so inconsistency follows. This establishes no converse:
sound deletion cannot by itself guarantee that separately compatible local
choices extend to one joint assignment. Allen's original section~4.3
explicitly reports this global limitation. The complete endpoint search
in the preceding proof is a different procedure.

For the body's containment example, \(I\,d\,J\) supplies
\(i^+<j^+\), while \(J\,b\,K\) supplies \(j^+<k^-\).
Transitivity gives \(i^+<k^-\), exactly \(I\,b\,K\).
The three numerical counterconfigurations for two intervals inside
\((0,10)\) satisfy their displayed endpoint inequalities directly.
For the switch/light/room calculation, composition distributes over
relation-set union: the initial \(L,R\) possibilities are
\(\{b,bi,o,m,di,s,si,fi,e\}\). Intersecting with
\(\{o,s,d\}\) leaves \(\{o,s\}\). The four relevant next
compositions are \(o\diamond o=\{b,o,m\}\),
\(o\diamond s=\{o\}\), \(m\diamond o=\{b\}\), and
\(m\diamond s=\{m\}\). Their union is \(\{b,o,m\}\),
whose intersection with \(\{b,m,mi,bi\}\) is \(\{b,m\}\).
Each composition can be checked by the endpoint partition above; no
numerical duration assumption enters this calculation.
\end{proof}
\end{annexitem}

\begin{annexitem}{annex:p19.rat132.cars}{thm:p19-rat132-cars}{Exhausting the car orders and their local transformations}
\begin{proof}
Write \(F,P,B,H,D\) for Ferrari, Porsche, Beetle, Hummer and Dodge.
The premises are \(F<P\), \(P<B\), \(P<H\), \(H<D\).
Transitivity fixes the chain \(F<P<H<D\). Since all positions are
distinct and \(B>P\), the Beetle can occur only between \(P,H\),
between \(H,D\), or after \(D\). Each choice respects every premise,
giving exactly \(FPBHD,FPHBD,FPHDB\).
All contain \(P<D\). The first two contain \(B<D\), while the
third contains \(D<B\), proving the stated necessary/possible claims.

An adjacent swap cannot alter the fixed chain without violating a premise.
Within the three solutions, swapping \(B,H\) takes the first to the
second, and swapping \(B,D\) takes the second to the third. Their
converses return along the path. The first and third differ in Beetle's
relative order with two cars, which one adjacent swap cannot change.
Thus their distance is exactly two in this move graph. This proves a
discrete graph fact, not a unique human construction sequence.
\end{proof}
\end{annexitem}

\begin{annexitem}{annex:p19.rat132.durative}{thm:p19-rat132-durative}{A preferred timeline, a consistent alternative and a contradiction}
\begin{proof}
Let \(A,F,S\) be alarm, fire and siren. With strict containment,
\(A\,d\,F\) means
\(f^-<a^-<a^+<f^+\), and \(S\,b\,A\) means
\(s^+<a^-\). The assignment
\(F=(0,8), A=(4,6), S=(-3,-1)\) satisfies these two premises
with siren before the fire. The assignment
\(F=(0,8), A=(4,6), S=(1,3)\) also satisfies them and additionally
satisfies \(S\,d\,F\), witnessing the consistency of the
three-premise description with siren during fire. The first assignment
is not a counterexample to that description's consistency: one failed
candidate is not a refutation of existence.

If instead \(F\,d\,S\), then \(f^+<s^+\). Combining the
premises yields
\(a^-<a^+<f^+<s^+<a^-\), impossible in a strict real order.
Thus the alternative third premise makes the set inconsistent.
These conclusions use proper intervals and strict enclosure. Reinterpreting
during as any overlap changes the relation semantics and therefore may
change the label of the inconsistent item; it does not invalidate the
explicit witness for the consistent enclosure item.
\end{proof}
\end{annexitem}
